% ============================================================
%  LAMPIRAN D – TEMPLAT PROMPT
%  Compile standalone:  pdflatex lampiran-d.tex
%  Compile as part of thesis: pdflatex main.tex
% ============================================================
\documentclass[../../thesis]{subfiles}
\usepackage{hyphenat}
\usepackage{iftex}

% ── Encoding, Language & Fonts ───────────────────────────────
\ifXeTeX
    \usepackage{polyglossia}
    \setmainlanguage{bahasai}
    \usepackage{fontspec}
    \setmainfont{TeX Gyre Termes}
    \setsansfont{TeX Gyre Heros}
    \setmonofont[Scale=0.88]{TeX Gyre Cursor}
\else
    \usepackage[utf8]{inputenc}
    \usepackage[T1]{fontenc}
    \usepackage[indonesian]{babel}
    \usepackage{mathptmx}
    \usepackage{helvet}
    \usepackage{courier}
\fi

% ── Page geometry (A4, UI thesis margins) ────────────────────
\usepackage[
  a4paper,
  top    = 2.54cm,
  bottom = 2.54cm,
  left   = 3.17cm,
  right  = 2.54cm
]{geometry}

% ── Line spacing & paragraph ─────────────────────────────────
\usepackage{setspace}
\onehalfspacing
\setlength{\parindent}{1.27cm}
\setlength{\parskip}{0pt}

% ── Microtype & justification ────────────────────────────────
\usepackage{microtype}
\sloppy

% ── Headings ─────────────────────────────────────────────────
\usepackage{titlesec}

\titleformat{\chapter}[block]
  {\centering\sffamily\bfseries\large}{}
  {0pt}{\MakeUppercase}
\titlespacing*{\chapter}{0pt}{0pt}{1.5em}

\titleformat{\section}
  {\sffamily\bfseries\normalsize}{\thesection}{1em}{}
\titlespacing*{\section}{0pt}{1.2em}{0.6em}

% ── Code listings (only used when compiled standalone; inside
%    main.tex the preamble from main.tex is used instead) ──────
\usepackage{listings}
\lstdefinestyle{promptbox}{
  basicstyle      = \ttfamily\scriptsize,
  breaklines      = true,
  breakatwhitespace = false,
  columns         = fullflexible,
  keepspaces      = true,
  showstringspaces = false,
  frame           = single,
  framesep        = 5pt,
  xleftmargin     = 10pt,
  xrightmargin    = 2pt,
  aboveskip       = 8pt,
  belowskip       = 8pt,
}

% ── Hyperlinks ───────────────────────────────────────────────
\usepackage[hidelinks]{hyperref}

% ============================================================
\begin{document}

% When standalone: switch to appendix lettering and set the counter
% so \chapter{} renders as "LAMPIRAN D". When compiled inside main.tex,
% main.tex issues \appendix and advances the counter before \subfile.
\ifSubfilesClassLoaded{}{\appendix\setcounter{chapter}{3}}

\chapter{Templat Prompt}
\label{lamp:prompt}

Lampiran ini memuat templat \textit{prompt} lengkap yang digunakan pada
\textit{pipeline} implementasi (Bab~4), yakni: (1)~\textit{prompt} ekstraksi
\textit{knowledge graph} dari teks buku (Pass-1), (2)~\textit{prompt}
klasifikasi relasi lintas-buku pada tahap \textit{completion} (Pass-2), dan
(3)~\textit{prompt} \textit{LLM-as-a-judge} untuk penyelesaian ketidaksepakatan
pakar pada konsensus Fase~1. Seluruh \textit{prompt} ditranskripsi apa adanya
dari kode implementasi final (\textit{notebook} \texttt{TA\_KG\_PIPELINE.ipynb}
untuk Pass-1 dan Pass-2, serta \texttt{KG\_CONSENSUS.ipynb} untuk
\textit{LLM-as-a-judge}). Kunci API dan kredensial basis data yang ada pada
kode sengaja \textbf{tidak} disertakan. Bagian bertanda kurung kurawal
(mis.\ \texttt{\{grade\}}, \texttt{\{text\_chunk\}}) merupakan \textit{placeholder}
yang disubstitusi pada saat eksekusi.

% ------------------------------------------------------------------
\section{Prompt Ekstraksi Pengetahuan (Pass-1)}
\label{lamp:prompt-ekstraksi}

\textit{Prompt} ekstraksi dikirim per potongan teks (\textit{chunk};
\texttt{CHUNK\_SIZE}~$=800$, \texttt{CHUNK\_OVERLAP}~$=200$) ke model
\texttt{gemini-2.5-flash}. \textit{Prompt} disusun dari delapan blok: (1)~instruksi
peran dan hierarki, (2)~konteks bab, (3)~daftar subbab, (4)~materi pokok dari
buku guru, (5)~glosarium, (6)~daftar konsep dari buku lain (kandidat tautan
lintas-buku), (7)~aturan penautan lintas-buku, dan (8)~spesifikasi format
keluaran JSON yang ketat (\textit{schema-constrained}). Blok 3--6 hanya
disisipkan jika data terkait tersedia. Templat utama ditunjukkan pada
Kode~\ref{lst:prompt-ekstraksi}, sedangkan keempat blok dinamis pada
Kode~\ref{lst:prompt-ekstraksi-blok}.

\begin{lstlisting}[style=promptbox,caption={Templat utama prompt ekstraksi (Pass-1).},label={lst:prompt-ekstraksi}]
Kamu membangun knowledge graph dari 3 buku: Biologi, Fisika, Kimia Kelas XII.
HIERARKI: Kelas -> Bab -> Subchapter -> Konsep

Buku: {grade} | Bab: {chapter}
Bab sebelumnya: {previous_chapter} | Bab berikutnya: {next_chapter}

{BLOK SUBCHAPTER}
{BLOK MATERI POKOK}
{BLOK GLOSARIUM}
{BLOK KONSEP LINTAS-BUKU}

ATURAN CROSS-BOOK:
1. cross_book_links HARUS menghubungkan ke buku lain, BUKAN buku saat ini ({grade}).
2. target_book hanya boleh: {allowed_targets}
3. Isi 1-3 cross_book_links jika ada keterkaitan kuat, atau [] jika tidak ada.

OUTPUT FORMAT (STRICT JSON):
{
  "chapter_summary": "Ringkasan 3-5 kalimat.",
  "subtopics": [{
    "name": "...",
    "concepts": [{
      "name": "...", "description": "Deskripsi 1-2 kalimat.",
      "glossary_validated": true/false, "materi_pokok_ref": "...",
      "cross_book_links": [{
        "target_book": "...", "target_concept": "...",
        "relation_type": "SAMA_DENGAN|APLIKASI_DARI|PRASYARAT_UNTUK|MEMPERDALAM|BERKAITAN_DENGAN",
        "explanation": "..."
      }],
      "relations": [{
        "type": "MENDEFINISIKAN|MENYEBABKAN|MEMUNGKINKAN|MENGATUR|BAGIAN_DARI|BERINTERAKSI_DENGAN|BERGANTUNG_PADA|MEMPENGARUHI",
        "target": "...", "description": "..."
      }]
    }]
  }],
  "chapter_relations": [{
    "type": "PRASYARAT|MEMPERSIAPKAN", "target_chapter": "...",
    "description": "...", "concept_links": [{"source_concept": "...", "target_concept": "...", "explanation": "..."}]
  }]
}

Semua output Bahasa Indonesia.

Teks:
{text_chunk}
\end{lstlisting}

\begin{lstlisting}[style=promptbox,caption={Empat blok dinamis yang disisipkan ke templat utama bila datanya tersedia.},label={lst:prompt-ekstraksi-blok}]
# Blok 3 - SUBCHAPTER
SUBCHAPTER:
  {label}. {nama_subchapter}
Gunakan nama subchapter persis untuk field "name" di "subtopics".

# Blok 4 - MATERI POKOK (dari buku guru; maks. 25 butir)
MATERI POKOK (BUKU GURU):
  1. {materi_pokok}
  2. {materi_pokok}
  ...

# Blok 5 - GLOSARIUM (maks. 25 entri; definisi dipangkas 100 karakter)
GLOSARIUM:
  {istilah}: {definisi}
  ...

# Blok 6 - KONSEP DARI BUKU LAIN (maks. 20 konsep per buku; kandidat tautan lintas-buku)
KONSEP DARI {nama_buku_lain}:
  - {konsep}
  - {konsep}
  ...
\end{lstlisting}

% ------------------------------------------------------------------
\section{Prompt Klasifikasi Relasi Lintas-Buku (Pass-2)}
\label{lamp:prompt-completion}

Pada tahap \textit{completion}, setiap pasangan konsep kandidat (hasil pencarian
ANN dengan \texttt{COS\_THRESHOLD}~$=0{,}75$ dan \texttt{TOP\_K}~$=15$)
diklasifikasikan oleh \texttt{gemini-2.5-flash} (\texttt{temperature}~$=0{,}2$)
dengan keluaran JSON terstruktur. \textit{Prompt} terdiri atas \textit{system
prompt} (Kode~\ref{lst:prompt-sys}) yang memuat definisi kelima tipe relasi,
\textit{user prompt} (Kode~\ref{lst:prompt-user}) yang memuat deskripsi kedua
konsep beserta konteks graf, dan skema respons (Kode~\ref{lst:prompt-schema}).

\begin{lstlisting}[style=promptbox,caption={System prompt klasifikasi relasi (Pass-2). Blok \texttt{DEFINISI TIPE RELASI} disisipkan dari konstanta \texttt{TYPE\_DEFINITIONS}.},label={lst:prompt-sys}]
Anda adalah ahli kurikulum sains SMA yang menentukan jenis keterkaitan
lintas-buku antara dua konsep dari Mata Pelajaran berbeda (Biologi/Fisika/Kimia Kelas XII).

Pilih SATU dari 6 nilai untuk relation_type (5 tipe LINTAS_BUKU_* + NONE jika tidak ada keterkaitan).

DEFINISI TIPE RELASI:

1. LINTAS_BUKU_SAMA_DENGAN     -- A dan B merujuk pada entitas/konsep yang sama di buku berbeda. Simetris.
2. LINTAS_BUKU_APLIKASI_DARI   -- A adalah penerapan/manifestasi konsep B di disiplin lain. Asimetris.
3. LINTAS_BUKU_PRASYARAT_UNTUK -- A adalah prasyarat untuk memahami B di buku lain. Asimetris, hierarkis.
4. LINTAS_BUKU_MEMPERDALAM     -- A memperdalam/memperluas pemahaman konsep B di buku lain. Asimetris.
5. LINTAS_BUKU_BERKAITAN_DENGAN -- Berkaitan secara umum lintas-buku (fallback). Simetris.

ATURAN:
- Hanya LINTAS_BUKU_* jika Mata Pelajaran A != B. Jika sama, kembalikan NONE.
- Pilih NONE jika tidak ada relasi jelas.
- Tipe asimetris: tentukan direction A_TO_B atau B_TO_A.
- SAMA_DENGAN dan BERKAITAN_DENGAN selalu SYMMETRIC.
- evidence_quote: kutip dari deskripsi konsep (bukan KONTEKS GRAF).
- confidence: 0.0-1.0, gunakan 0.85+ hanya jika sangat yakin.
\end{lstlisting}

\begin{lstlisting}[style=promptbox,caption={User prompt klasifikasi relasi (Pass-2). Blok \texttt{KONTEKS GRAF} memuat subtopik, hingga lima edge keluar tiap konsep, dan tetangga bersama.},label={lst:prompt-user}]
Tentukan relasi LINTAS_BUKU antara dua konsep.

KONSEP A:  Buku={a.book} | Bab={a.chapter} | Subtopic={a.subtopic}
  Nama: {a.name}
  Deskripsi: {a.description}

KONSEP B:  Buku={b.book} | Bab={b.chapter} | Subtopic={b.subtopic}
  Nama: {b.name}
  Deskripsi: {b.description}

Embedding similarity: {similarity}

KONTEKS GRAF (HUBUNGAN YANG SUDAH ADA)
Konsep A ("{a.name}") Subtopic: {a.subtopic}
Edges A:
  - {relasi} -> "{target}"
  ... +{n} edge lain

Konsep B ("{b.name}") Subtopic: {b.subtopic}
Edges B:
  - {relasi} -> "{target}"
  ... +{n} edge lain

Tetangga bersama: {daftar_tetangga_bersama}

Output JSON sesuai schema.
\end{lstlisting}

\begin{lstlisting}[style=promptbox,caption={Skema respons JSON yang dipaksakan pada keluaran model (Pass-2).},label={lst:prompt-schema}]
{
  "type": "object",
  "properties": {
    "relation_type":  {"type": "string",
                        "enum": [<5 tipe LINTAS_BUKU_*>, "NONE"]},
    "direction":      {"type": "string",
                        "enum": ["A_TO_B", "B_TO_A", "SYMMETRIC"]},
    "confidence":     {"type": "number"},
    "evidence_quote": {"type": "string"},
    "rationale":      {"type": "string"}
  },
  "required": ["relation_type", "direction", "confidence",
               "evidence_quote", "rationale"]
}
\end{lstlisting}

\noindent Prediksi diterima hanya bila \texttt{relation\_type} bukan
\texttt{NONE}, \texttt{confidence} melampaui ambang per-tipe
(\texttt{LINTAS\_BUKU\_BERKAITAN\_DENGAN}~$\geq 0{,}85$; tipe lain
$\geq 0{,}70$), dan arah konsisten dengan sifat simetri tipe relasi.

% ------------------------------------------------------------------
\section{Prompt LLM-as-a-Judge (Konsensus Fase~1)}
\label{lamp:prompt-judge}

Pada konsensus Fase~1, kasus ketidaksepakatan antara dua pakar diselesaikan oleh
model \textit{LLM-as-a-judge} (model Gemini Pro, dikonfigurasi melalui variabel
lingkungan). \textit{Judge} menerima konteks pendukung berupa potongan teks
\textit{e-book} yang relevan (\textit{retrieval} sederhana berbasis kecocokan
kata kunci) dan mengembalikan keputusan dalam format JSON dengan kunci
\texttt{reasoning} dan \texttt{final\_consensus\_label}. Terdapat dua varian:
penilaian ketidaksepakatan label (Kode~\ref{lst:judge-disagreed}) dan penilaian
\textit{missing triple} yang hanya diusulkan satu pakar (Kode~\ref{lst:judge-missing}).

\begin{lstlisting}[style=promptbox,caption={Prompt judge untuk ketidaksepakatan label rating.},label={lst:judge-disagreed}]
You are an expert KG validator for high-school chemistry curriculum.
Decide final label for this disagreement:
- triple_id: {triple_id}
- chapter: {chapter}
- subject: {subject}
- relation: {relation}
- target: {target}
- description: {description}
- reviewer_1_label: {r1_rating}
- reviewer_2_label: {r2_rating}

Relevant ebook context:
{ebook_context}
Use the ebook context as supporting evidence when it is relevant. Do not invent facts beyond the triple and context.
Return strict JSON with keys: reasoning, final_consensus_label.
Use one of these labels: correct, wrong, partial, missing, abaikan.
If final_consensus_label is partial or missing, write reasoning in Bahasa Indonesia as a clean KG relation description that can replace or enrich the original description.
For other labels, write a concise validation rationale in Bahasa Indonesia.
\end{lstlisting}

\begin{lstlisting}[style=promptbox,caption={Prompt judge untuk missing triple yang diusulkan satu pakar.},label={lst:judge-missing}]
You are validating a missing triple proposed by only one reviewer.
Decide if this triple should be accepted as curriculum-required 'missing' or ignored as 'Abaikan'.
- reviewer_source: {reviewer}
- chapter: {chapter}
- subject: {subject}
- relation: {relation}
- target: {target}
- description: {description}

Relevant ebook context:
{ebook_context}
Use the ebook context as supporting evidence when it is relevant. Do not invent facts beyond the triple and context.
Return strict JSON with keys: reasoning, final_consensus_label.
Use one of these labels: missing, abaikan.
If final_consensus_label is missing, write reasoning in Bahasa Indonesia as a clean KG relation description that can be used directly as the description.
If final_consensus_label is abaikan, write a concise validation rationale in Bahasa Indonesia.
\end{lstlisting}

\end{document}
